\documentclass[a1paper,portrait]{bxjsarticle}
\usepackage[dvipdfmx]{graphicx}
\usepackage[dvipdfmx]{color}
\usepackage{amsmath}
\usepackage{multicol}
\usepackage{enumitem}
\usepackage{bm}
\usepackage{geometry}
\usepackage{bicaption}
\usepackage{algorithmic}
\geometry{margin=2cm,top=1cm, bottom=1cm}
\usepackage{tikz}
\usepackage{tcolorbox}
\tcbuselibrary{breakable,skins}
\usepackage{caption}
\usepackage{subcaption}
\usepackage{pgfplots}
\usepackage{float}
\usepackage{comment}
\usepackage{anyfontsize}
\usepackage{xcolor}
\usetikzlibrary{positioning}
\usetikzlibrary{calc}
\pgfplotsset{compat=1.16}
\usepackage{xcolor} % プリアンブルに追加
\definecolor{comsolblue}{RGB}{40,40,120}
\definecolor{comsolpink}{RGB}{180,40,100}

% Times New Roman フォント設定
% Use Times for math symbols, but set default text font to sans-serif for template matching.
\usepackage{newtxtext,newtxmath}
\renewcommand{\familydefault}{\sfdefault}

\usepackage{scalefnt}

% 日本語と英数字のフォントサイズ統一のためのコマンド定義
\newcommand{\timesroman}[1]{\fontfamily{ptm}\selectfont #1}

% 色の定義（PDFの配色を参考）
\definecolor{headercolor}{RGB}{0,51,153}
\definecolor{sectionblue}{RGB}{51,102,204}
\definecolor{sectiongreen}{RGB}{102,153,51}
\definecolor{boxcolor}{RGB}{240,245,255}
\definecolor{resultbox}{RGB}{255,250,240}
\definecolor{lightsectionblue}{RGB}{200,220,255}
\definecolor{lightyellow}{RGB}{255, 255, 204}

\makeatletter
\DeclareMathSizes{40}{40}{32}{24}
\makeatother
\DeclareCaptionFont{biglabel}{\fontsize{30pt}{36pt}\selectfont}
\captionsetup[figure]{labelfont=biglabel, textfont=normalsize}

% セクションのスタイル設定
\usepackage{titlesec}
\titleformat{\section}
{\normalfont\Huge\bfseries\color{white}}
{\Roman{section}.}{0.5em}{}
\titlespacing*{\section}{0pt}{0pt}{10pt}

% サブセクションのスタイル
\titleformat{\subsection}
{\normalfont\Large\bfseries\color{sectionblue}}
{\thesubsection}{0.5em}{}

% tcolorboxの設定
\tcbset{
sectionbox/.style={colback=sectionblue,coltext=white,
boxrule=0pt,
arc=5mm,
left=15pt,
right=15pt,
top=10pt,
bottom=10pt,
fontupper=\Large\bfseries
},
subsectionbox/.style={colback=sectionblue!15,   % 薄い水色背景
coltext=sectionblue!80!black, % 文字はやや濃い青
boxrule=0pt,
arc=3mm,
left=12pt,right=12pt,top=6pt,bottom=6pt,
fontupper=\large\bfseries
},
contentbox/.style={colback=white,colframe=sectionblue,
boxrule=2pt,
arc=3mm,
boxsep=10pt,
left=10pt,
right=10pt,
top=10pt,
bottom=10pt
},
subcontentbox/.style={colback=white,           % 背景は白
colframe=sectionblue,   % 枠を薄い黄色
boxrule=2pt,             % 枠の太さ
arc=3mm,                 % 角を丸く
boxsep=6pt,              % 中身と枠の余白
left=5pt,right=5pt,top=1pt,bottom=1pt,
enhanced,                % 高度なオプション有効
},
resultbox/.style={colback=resultbox,colframe=sectiongreen,
boxrule=2pt,
arc=3mm,
boxsep=10pt
}
}

\begin{document}

% ヘッダー部分
\begin{flushleft}
{\fontsize{60}{72}\selectfont\bfseries\color{headercolor} 電磁場開領域問題におけるケルビン変換}

\vspace{5pt}

{\fontsize{40}{48}\selectfont\bfseries\color{headercolor} \quad ~~Maxwell方程式の共形変換を用いた外部領域の定式化}

\vspace{5pt}

{\fontsize{35}{42}\selectfont 近畿大学 理工学部 電気電子工学科 CAE-AI研究室 \quad 菅原賢悟}

\vspace{5pt}

\hrule
\end{flushleft}


\vspace{10pt}

% I. Introduction（全幅で配置）
\begin{tcolorbox}[sectionbox]
\begin{center}
{\fontsize{36pt}{40pt}\selectfont はじめに}
\end{center}
\end{tcolorbox}
{\fontsize{24pt}{33pt}\selectfont\par
ケルビン変換は無限空間を有限空間へ写像する座標変換であり、Maxwell方程式の共形変換として機能する。本研究では、微分幾何学の概念を適用し、外部領域における材料定数の計量および空間依存性を導出した。これにより、外部領域にも物体や材料特性、PMLを含めることが可能となる。低周波の渦電流問題と高周波の放射問題の両方に適用可能である。COMSOL Multiphysics®による数値計算例で本手法の妥当性を検証した。
\par}

\vspace{5pt}

\setlength{\columnseprule}{0pt}
\begin{multicols}{2}

\begin{tcolorbox}[sectionbox]
\begin{center}
{\fontsize{36pt}{43pt}\selectfont ケルビン変換の定式化と電磁場問題への応用}
\end{center}
\end{tcolorbox}

	\begin{tcolorbox}[contentbox,breakable]
		\begin{tcolorbox}[subsectionbox]
		{\fontsize{34pt}{40pt}
			\selectfont 従来提案されているケルビン変換
		}
		\end{tcolorbox}

		\begin{center}
		\includegraphics[clip, width=0.65\linewidth]{figures/Kelvin_Concept.pdf}
		\captionof{figure}{\fontsize{26pt}{30pt}\selectfont ケルビン変換の概念}
		\end{center}

		\vspace{10pt}

		{\fontsize{20pt}{26pt}\selectfont
		\begin{enumerate}[label={[\arabic*]}, leftmargin=*, itemsep=0pt]
		\item E.~M.~Freeman, D.~A.~Lowther, "A novel mapping technique for open boundary finite element solutions to Poisson's equation," IEEE Trans. Magn., vol.24, no.6, pp.2934-2936, 1988.
		\end{enumerate}
		}
	\end{tcolorbox}

	\begin{tcolorbox}[contentbox,breakable]
		\begin{tcolorbox}[subsectionbox]
		{\fontsize{34pt}{40pt}\selectfont 
		拡張型ケルビン変換の定式化
		}
		\end{tcolorbox}
		\begin{center}
		\includegraphics[clip, width=0.65\linewidth]{figures/formulation.pdf}
		\captionof{figure}{\fontsize{26pt}{30pt}\selectfont Maxwell方程式の共形変換}
		\end{center}
		{\fontsize{24pt}{33pt}\selectfont\par
		Maxwell方程式の共形変換により、外部領域の材料定数（透磁率・誘電率・導電率）の空間依存性を計量から導出。2次元・軸対称では$(a/r)^4$、3次元では$(a/r)^2$で変調される。COMSOL Multiphysics®の係数型PDEインターフェースで実装。
		\par}
	\end{tcolorbox}

	\begin{tcolorbox}[contentbox,breakable]
		\begin{tcolorbox}[subsectionbox]
		{\fontsize{34pt}{40pt}\selectfont 想定される応用問題}
		\end{tcolorbox}
		\begin{center}
		\includegraphics[clip, width=0.65\linewidth]{figures/Kelvin_application.pdf}
		\captionof{figure}{\fontsize{26pt}{30pt}\selectfont 外部領域に散乱体を含む電磁場の放射問題}
		\end{center}
		{\fontsize{24pt}{33pt}\selectfont\par
		外部領域にも物体や材料が存在する場合にも、ケルビン変換により無限領域を有限領域に写像することで、効率的な解析が可能。
		\par}
	\end{tcolorbox}

\columnbreak

\begin{tcolorbox}[sectionbox]
	\begin{center}
	{\fontsize{36pt}{43pt}\selectfont 拡張型ケルビン変換の適用事例}
	\end{center}
	\end{tcolorbox}

	\begin{tcolorbox}[contentbox,breakable]
		\begin{tcolorbox}[subsectionbox]
		{\fontsize{34pt}{40pt}\selectfont 事例A：低周波渦電流探傷問題 }
		\end{tcolorbox}

		\begin{center}
		\includegraphics[clip, width=0.50\linewidth]{figures/3d-example-ECT_earth.pdf}\hfill
		\includegraphics[clip, width=0.50\linewidth]{figures/3d-example-kelvin-transform.pdf}
		\captionof{figure}{\fontsize{26pt}{30pt}\selectfont 渦電流探傷解析}
		\end{center}

		{\fontsize{24pt}{33pt}\selectfont\par
		外部境界に導体（軟磁性材料）を含めた渦電流探傷解析を、COMSOLのAC/DCモジュールで実施。複雑形状導体の渦電流分布を高精度・高速に解析できることを確認。
		\par}

		\vspace{10pt}

		{\fontsize{20pt}{26pt}\selectfont
		\begin{enumerate}[label={[\arabic*]}, leftmargin=*, itemsep=0pt, start=2]
		\item K. Sugahara, "Electromagnetic Analysis of Eddy Current Testing With Kelvin Transformation," IEEE Trans. Magn., vol.58, no.9, pp.1-6, Sept. 2022.
		\end{enumerate}
		}
	\end{tcolorbox}

	\begin{tcolorbox}[contentbox,breakable]
		\begin{tcolorbox}[subsectionbox]
		{\fontsize{34pt}{40pt}\selectfont
		事例B：高周波放射散乱問題（アンテナ解析）
		}
		\end{tcolorbox}

		\begin{center}
		\begin{minipage}{0.45\linewidth}
			\centering
			\includegraphics[clip, width=0.9\linewidth]{figures/HertzianDipole.pdf}
			\captionof{figure}{\fontsize{26pt}{30pt}\selectfont 電磁場散乱問題}
		\end{minipage}\hfill
		\begin{minipage}{0.55\linewidth}
			\centering
			\includegraphics[clip, width=\linewidth]{figures/Figure_scatter_xy_E_abs.pdf}
			\captionof{figure}{\fontsize{30pt}{36pt}\selectfont 電界分布の比較}
		\end{minipage}
		\end{center}

		{\fontsize{24pt}{33pt}\selectfont\par
		ヘルツアンテナの放射問題に適用し、解析解との高精度な一致を確認。PMLを外部領域に配置することで、従来のPMLやインピーダンス境界条件を用いず効率的な放射界解析を実現。
		\par}

		\vspace{10pt}

		{\fontsize{20pt}{26pt}\selectfont
		\begin{enumerate}[label={[\arabic*]}, leftmargin=*, itemsep=0pt, start=3]
		\item K. Sugahara, "Extended Kelvin Transformation for Solving Radiating Electromagnetic Fields," IEICE Trans. Electron., vol.E108-C, no.4, pp.189-194, April 2025.
		\end{enumerate}
		}
	\end{tcolorbox}

\begin{tcolorbox}[sectionbox]
	\begin{center}
	{\fontsize{36pt}{43pt}\selectfont まとめ}
	\end{center}
	\end{tcolorbox}

	\begin{tcolorbox}[contentbox,breakable]
	{\fontsize{24pt}{33pt}\selectfont\par
	Maxwell方程式の共形変換を用いた拡張ケルビン変換により、低周波から高周波まで単一手法で解析可能なことをCOMSOL Multiphysics®を用いて実証した。
	\par
	\begin{itemize}
		\item 周波数帯域ごとの手法切替が不要
		\item 外部境界に材料特性・PMLを含められる
		\item 広帯域電磁界現象の統一的解析を実現
	\end{itemize}
	}\end{tcolorbox}
\end{multicols}

\begin{tikzpicture}[remember picture,overlay]
\shade[left color=comsolblue, right color=comsolpink] 
	([yshift=0pt]current page.south west) rectangle 
	([yshift=40pt]current page.south east);

\node[anchor=south, yshift=0pt, white] at (current page.south) {
	\fontsize{24}{40}\selectfont Excerpt from proceedings of the 
	\textbf{COMSOL CONFERENCE} 2025 TOKYO
};
\end{tikzpicture}

\end{document}
